\documentclass[11pt,a4paper]{article}
\usepackage{amsmath,amssymb,amsthm}
\usepackage[margin=2.5cm]{geometry}
\usepackage{hyperref}

\newtheorem{definition}{Definition}[section]
\newtheorem{theorem}[definition]{Theorem}
\newtheorem{proposition}[definition]{Proposition}
\newtheorem{remark}[definition]{Remark}
\newtheorem{conjecture}[definition]{Conjecture}

\title{Hyperseed Formalization:\\
  What Is Science, Exactly?}
\author{ProtomegaTron (automated formalization)}
\date{2026-09-18}

\begin{document}
\maketitle

\begin{abstract}
We formalize the intellectual structure of Ben Goertzel's Substack article
``What Is Science, Exactly?'' (2026-01-02) within the Hyperseed ontology.
The article examines science as a social-epistemic process rather than a
single method. Key mappings: science as communal fiber construction,
reproducibility as fiber verification, Bayesian updating as fiber revision,
reproducibility crisis as fiber-base decoupling, peer review as inter-observer
fiber comparison, and AI science tools as base interface expansion.
\end{abstract}

\tableofcontents

%% =================================================================
\section{Communal fiber construction: science as social process}
\label{sec:communal}
%% =================================================================

\begin{theorem}[Science as communal fiber --- claims 1--9, 25]
\label{thm:communal}
There is no single ``scientific method.'' Science is a family of practices
organized around structured skepticism, reproducibility, and communal
evidence evaluation. The textbook version (hypothesis $\to$ experiment
$\to$ theory) is a pedagogical simplification.

Science is fundamentally social: individual brilliance matters, but the
epistemic power comes from the community process --- peer review,
replication, critique, synthesis.

In Hyperseed terms, science is \emph{communal fiber construction}: many
observers contributing to a shared fiber (theory) over a shared base
(evidence):
\[
  F_{\text{science}} = \bigcup_{i \in \text{community}}
  f_i(B_{\text{shared}})
\]

No single observer builds the full fiber. Each contributes observations,
interpretations, and critiques that collectively shape the community
fiber. The result is community-wide knowledge --- not what any individual
knows but what the community collectively knows through its distributed
evidence evaluation process.

Key mechanisms of communal fiber construction:
\begin{itemize}
\item \textbf{Peer review} = inter-observer fiber comparison: different
  observers compare their fibers (interpretations) over the same base
  (data/methods), and discrepancies signal errors or ambiguities.
\item \textbf{Replication} = fiber verification: checking that the same
  base produces the same fiber when traversed by different observers.
\item \textbf{Critique} = fiber error detection: identifying where the
  community fiber doesn't faithfully represent the base.
\item \textbf{Synthesis} = fiber integration: combining partial fibers
  from different observers into a coherent whole.
\end{itemize}

The demarcation problem (what distinguishes science from non-science)
maps to the question of what distinguishes rigorous communal fiber
construction from other kinds of collective belief formation. No clean
criterion captures all of science --- the boundary is fuzzy because
the fiber construction practices vary across domains.
\end{theorem}

%% =================================================================
\section{Fiber revision: the Bayesian core}
\label{sec:bayesian}
%% =================================================================

\begin{proposition}[Bayesian fiber revision --- claims 10--13, 27]
\label{prop:bayesian}
At its best, science approximates Bayesian updating on a community-wide
prior. Each experiment revises the community fiber based on new base
evidence:
\[
  F_{t+1} = \text{Update}(F_t, e_{\text{new}}, w(e_{\text{new}}))
\]

where $w(e)$ is the evidential weight: controlled experiments weigh
more than observational studies, large samples more than small ones,
pre-registered studies more than post-hoc analyses.

The community prior evolves through accumulation of evidence, theory
development, and paradigm shifts. In Hyperseed terms, this is the
fiber's temporal evolution: the community fiber at time $t+1$ depends
on the fiber at time $t$ plus new base evidence.

Real science deviates from ideal Bayesianism in systematic ways:
confirmation bias (over-weighting evidence that confirms the current
fiber), anchoring (the prior fiber has too much inertia), social
pressure (the fiber reflects community politics, not just evidence),
and career incentives (the fiber reflects what's rewarded, not just
what's true).

These deviations are \emph{fiber distortions}: systematic ways in which
the community fiber fails to faithfully represent the base evidence.
Understanding these distortions is essential for improving the
fiber-construction process.
\end{proposition}

%% =================================================================
\section{Fiber-base decoupling: the reproducibility crisis}
\label{sec:crisis}
%% =================================================================

\begin{theorem}[Reproducibility crisis --- claims 14--18, 28]
\label{thm:crisis}
The current reproducibility crisis reveals what happens when the
social-epistemic machinery breaks down --- when the fiber (published
claims) decouples from the base (actual evidence):
\[
  \|F_{\text{published}} - F_{\text{evidence}}\| \gg 0
  \quad \text{(fiber-base decoupling)}
\]

The decoupling is driven by incentive misalignment:
\begin{itemize}
\item \textbf{Publish-or-perish}: rewards fiber quantity over fiber
  quality. The system incentivizes producing more fiber (papers)
  regardless of whether it faithfully represents the base (evidence).
\item \textbf{Publication bias}: journals preferentially accept fiber
  that shows positive results, creating a distorted community fiber.
  The ``file drawer problem'' means negative-result fiber disappears.
\item \textbf{Statistical malpractice}: p-hacking, HARKing, underpowered
  studies --- systematic methods for making the fiber \emph{appear} to
  represent the base when it doesn't.
\end{itemize}

The crisis extends beyond psychology to biomedical research, economics,
and other fields. It's a systemic fiber-base decoupling, not a
disciplinary one.

In Hyperseed terms, the reproducibility crisis is a failure of
\emph{fiber verification}: the community has stopped systematically
checking that the fiber faithfully represents the base. Without
verification (replication), the fiber floats free of its evidential
base --- claims accumulate without adequate base support.
\end{theorem}

%% =================================================================
\section{AI as base interface expansion for science}
\label{sec:ai}
%% =================================================================

\begin{proposition}[AI and science --- claims 19--24, 30]
\label{prop:ai}
AI could either help fix science or make it worse. In Hyperseed terms,
AI tools for science expand the base interface (faster analysis, more
data processing) but don't necessarily improve the fiber (better
theories, better reasoning):

\textbf{Helpful applications (base interface improvement):}
\begin{itemize}
\item Better statistical analysis (detecting fiber-base discrepancies).
\item Automated replication checking (systematic fiber verification).
\item Bias detection (identifying fiber distortions).
\item Literature synthesis (integrating partial fibers).
\end{itemize}

\textbf{Harmful applications (base interface abuse):}
\begin{itemize}
\item P-hacking at scale (automated fiber-base decoupling).
\item Publication mills (mass-produced unreliable fiber).
\item Deepfake data (corrupting the base itself).
\end{itemize}

AI is an amplifier: it amplifies whatever the system incentivizes. If
incentives are broken, AI amplifies the brokenness. This mirrors the
OpenClaw analysis: better hands for a scientific brain. The hands
(AI tools) amplify what the brain (scientific community) does, for
better or worse.

AGI-as-scientist requires genuine reasoning capabilities (abstraction,
causal inference, creative hypothesis generation) AND social integration
(participating in the communal epistemic process). This requires both
fiber improvement (better reasoning) and base integration (participating
in the community fiber construction).
\end{proposition}

%% =================================================================
\section{Cross-references}
%% =================================================================

\begin{remark}[Connections to other formalizations]
\begin{itemize}
\item \textbf{OpenClaw: Amazing Hands for a Brain} (2026-02-03): AI tools
  as base interface expansion. The same hands-vs-brain distinction applies
  to AI-for-science.
\item \textbf{Evidence Is to Logic What Energy Is to Physics} (2026-03-10):
  evidence conservation. The anti-hallucination theorems apply to science:
  scientific claims should be bounded by evidence (hallucination bound).
\item \textbf{Goertzel vs Epstein} (2026-02-20): provenance. Scientific
  provenance (who found what, funded by whom) is part of the fiber's
  evidence trail.
\item \textbf{LeCun's SAI Is a Special Case of AGI} (2026-03-07): credit.
  Credit erasure in science is fiber provenance deletion --- the same
  pattern LeCun exhibits with AGI research.
\end{itemize}
\end{remark}

\begin{conjecture}[Fiber verification as science's defining property]
Conjecture: the defining property of science (what distinguishes it from
other forms of collective belief) is systematic fiber verification ---
the commitment to checking that the communal fiber faithfully represents
the base evidence, through replication, peer review, and structured
skepticism.

If this conjecture holds, the reproducibility crisis is not a peripheral
problem but an existential threat to science itself: when fiber
verification breaks down, science ceases to be science and becomes
another form of collective myth-making. Restoring science requires
restoring fiber verification --- which means fixing the incentive
structures that currently discourage it.
\end{conjecture}

\end{document}
